\begin{abstract}
A problemática central deste trabalho consiste em investigar se modelos de linguagem de grande porte (LLMs), amplamente utilizados em domínios como medicina e educação, são capazes de atender às demandas específicas da pecuária leiteira, fornecendo respostas úteis, coerentes e contextualizadas para auxiliar produtores na tomada de decisão. 

Para isso, o estudo realiza uma análise comparativa entre diferentes modelos generalistas, avaliando seu desempenho no domínio agropecuário. Como abordagem metodológica, são empregadas métricas consolidadas de avaliação textual e semântica, incluindo BLEU e ROUGE, para similaridade lexical, bem como BERTScore e BARTScore, para análise de proximidade semântica. Adicionalmente, utiliza-se o GPTScore para incorporar uma perspectiva contextual qualitativa às avaliações.

O trabalho também explora estratégias de aprimoramento baseadas em Retrieval-Augmented Generation (RAG), incluindo a aplicação de variantes do método, parametrizações ajustáveis e a categorização das perguntas em diferentes seções temáticas do domínio. Essas abordagens visam melhorar a relevância das respostas e a aderência ao contexto específico da pecuária leiteira.

Os resultados indicam que, embora os modelos generalistas apresentem desempenho razoável, existem limitações importantes quando aplicados a domínios altamente especializados. A incorporação de técnicas como RAG, aliada a estratégias de organização e parametrização das consultas, demonstra potencial significativo para mitigar essas limitações e aprimorar a qualidade das respostas geradas.

\textit{Palavras-chave}: Large Language Models, Pecuária Leiteira, Retrieval-Augmented Generation, Avaliação de Modelos, Processamento de Linguagem Natural.
\end{abstract}